\section{Introduction}

We set out to demonstrate that fairness rankings of language models are more stable than adversarial robustness rankings across distribution shifts --- and found that the mechanism we proposed to explain this difference was wrong. Both fairness and adversarial robustness exhibit strikingly high rank stability (partial Spearman $\rho > 0.68$ for all pairs, zero rank reversals) after controlling for general capability. The adversarial design of robustness benchmarks does not disrupt cross-split model ordering. Yet this reversal is the finding: trustworthiness rankings, across both dimensions we studied, are highly stable. And that stability has direct implications for how we evaluate and select language models for deployment.

When practitioners select a language model for a safety-sensitive application --- healthcare, legal information, public-facing dialogue --- they rely on trustworthiness evaluations: fairness audits on benchmark datasets, adversarial robustness scores, alignment measurements. The implicit assumption is that a model ranked high on fairness or robustness \emph{in the evaluation condition} will remain trustworthy \emph{in deployment conditions} that differ from the benchmark. This assumption is widely made but has not been empirically validated.

The gap is real and consequential. Prior multi-model trustworthiness evaluations --- TrustLLM \citep{Sun2024TrustLLM}, DecodingTrust \citep{Wang2023DecodingTrust}, HELM \citep{Liang2022Holistic} --- report absolute performance scores for models on individual benchmarks. None asks whether in-distribution model rankings predict out-of-distribution model rankings \emph{across dimensions} as a primary research question. The field's closest methodological precedent is \citet{GeversAndDaelemans2026}, who apply rank correlation with capability control to commonsense benchmarks --- but not to trustworthiness dimensions with their distinctive latent-property vs.\ adversarial-design structure.

This leaves a specific empirical gap: do in-distribution trustworthiness rankings predict out-of-distribution trustworthiness rankings at scale, and does the answer differ by dimension?

We fill this gap with a focused analysis: partial Spearman $\rho$ (MMLU-controlled) between in-distribution and out-of-distribution model rankings across three trustworthiness benchmark pairs --- BBQ-Disambig$\rightarrow$BBQ-Ambig (fairness), ANLI R1$\rightarrow$R3 (adversarial robustness), and OOD robustness --- for $N=16$ LLMs from TrustLLM \citep{Sun2024TrustLLM}. The key methodological contribution is the capability control: raw Spearman $\rho$ is near-ceiling for all dimensions ($\approx 0.98$), masking dimension-specific structure. After partial correlation removing MMLU rank, the fairness-robustness differential emerges.

Our findings revise the theoretical picture on two fronts. First, fairness rank ordering is highly stable across BBQ context splits even after capability control ($\rho = 0.962$, $p \approx 0$, CI$=[0.90, 1.00]$, $N=16$). This is consistent with fairness failures encoding stable latent properties of model representations --- stereotypical associations formed during pretraining persist across the BBQ context shift. Second, and contrary to our original hypothesis, adversarial robustness benchmark construction does \emph{not} disrupt cross-split rank stability. Both ANLI R1$\rightarrow$R3 ($\rho = 0.684$, $p = 0.007$) and OOD robustness ($\rho = 0.868$, $p < 0.001$) show significantly positive partial $\rho$, with zero rank reversals across $N=13$ models. The adversarial disruption mechanism --- grounded in the design philosophy of AdvGLUE and ANLI and a GPT-3.5/4 anecdote from DecodingTrust \citep{Wang2023DecodingTrust} --- is falsified at scale.

The directional hypothesis that fairness stability exceeds robustness stability receives moderate support ($\Delta\rho = 0.192$, Fisher $z$ $p = 0.024$, $N=13$), though the preregistered effect size criterion ($\Delta\rho \geq 0.200$) was not met, falling short by 0.008. We report this as directional evidence, not a confirmed finding.

This work makes the following contributions:

\begin{enumerate}
\item \textbf{Empirical baseline.} First computation of partial Spearman $\rho$ (MMLU-controlled) between in-distribution and OOD trustworthiness benchmark model rankings across 16 LLMs, establishing that both fairness ($\rho = 0.962$) and adversarial robustness ($\rho = 0.684$--$0.868$) rankings are highly stable across distribution shifts.

\item \textbf{Methodological demonstration.} MMLU-controlled partial Spearman $\rho$ is a tractable, data-efficient operationalization of trustworthiness benchmark predictive validity using published scores only --- no new data collection required.

\item \textbf{Theoretical revision.} The adversarial disruption hypothesis is falsified. Both fairness and robustness are stable latent model properties. The directional fairness advantage has no confirmed mechanism, setting a new research agenda.
\end{enumerate}

The paper is organized as follows. Section~\ref{sec:related} situates our work within prior trustworthiness evaluation, adversarial robustness, and cross-benchmark predictive validity literature. Section~\ref{sec:methodology} describes the data, benchmark pairs, and partial correlation methodology. Section~\ref{sec:experiments} presents our experimental design and research questions. Section~\ref{sec:results} reports results. Section~\ref{sec:discussion} discusses implications and limitations. Section~\ref{sec:conclusion} concludes.
